\documentclass[12pt,a4paper]{ctexart}

% ============================================================
%                       宏包设置
% ============================================================

\usepackage{amsmath, amssymb, amsfonts}
\usepackage{geometry}
\usepackage{fancyhdr}
\usepackage{graphicx}
\usepackage{xcolor}
\usepackage{setspace}
\usepackage[most]{tcolorbox}
\usepackage{enumitem}

% ============================================================
%                 清新蓝绿色配色
% ============================================================

% 主色：清爽青绿色
\definecolor{SZURed}{HTML}{2F9E8F}

% 题目区域：浅薄荷绿
\definecolor{SZULightRed}{HTML}{EAF7F4}

% 答案区域：浅天空蓝
\definecolor{SZULightRose}{HTML}{EEF7FC}

% 边框：柔和蓝绿色
\definecolor{SZUSoftBorder}{HTML}{9CCFC8}

% 强调文字：沉稳蓝绿色
\definecolor{SZUTextRed}{HTML}{256B72}

% ============================================================
%                       页面设置
% ============================================================

\geometry{
    left=2.5cm,
    right=2.5cm,
    top=2.5cm,
    bottom=2.5cm
}

\setlength{\parindent}{0pt}
\setlength{\parskip}{0.6em}
\linespread{1.15}

% ============================================================
%                       课程信息
% ============================================================

\newcommand{\coursename}{概率论与数理统计}
\newcommand{\department}{深圳大学生物医学工程}
\newcommand{\homeworknumber}{作业 3}

% ============================================================
%                  学生信息 —— 请学生修改
% ============================================================

\newcommand{\studentname}{请填写姓名}
\newcommand{\studentid}{请填写学号}

% ============================================================
%                       页眉页脚
% ============================================================

\pagestyle{fancy}
\fancyhf{}

\lhead{\textcolor{SZUTextRed}{\coursename}}
\rhead{\textcolor{SZUTextRed}{\homeworknumber}}
\cfoot{\textcolor{SZUTextRed}{\thepage}}

\renewcommand{\headrulewidth}{0.4pt}
\renewcommand{\headrule}{
    \hbox to\headwidth{
        \color{SZUSoftBorder}\leaders\hrule height \headrulewidth\hfill
    }
}

\setlength{\headheight}{15pt}

% ============================================================
%                    题目编号与环境设置
% ============================================================

\newcounter{problem}

% -------------------------
% 题目环境
% -------------------------

\newenvironment{problem}[1][]
{
    \stepcounter{problem}
    \begin{tcolorbox}[
        title={题目 \theproblem\if&#1&\relax\else\quad[#1]\fi},
        colback=SZULightRed,
        colframe=SZURed,
        coltitle=white,
        fonttitle=\bfseries,
        boxrule=0.8pt,
        arc=2mm,
        left=3mm,
        right=3mm,
        top=2.5mm,
        bottom=2.5mm,
        breakable
    ]
}
{
    \end{tcolorbox}
}

% -------------------------
% 答案环境
% -------------------------

\newenvironment{answer}
{
    \begin{tcolorbox}[
        title={学生答案},
        colback=SZULightRose,
        colframe=SZUSoftBorder,
        coltitle=SZUTextRed,
        fonttitle=\bfseries,
        boxrule=0.7pt,
        arc=2mm,
        left=3mm,
        right=3mm,
        top=2.5mm,
        bottom=2.5mm,
        breakable
    ]
}
{
    \end{tcolorbox}
}

% ============================================================
%                         正文
% ============================================================

\begin{document}

% ============================================================
%                         标题
% ============================================================

\begin{center}

    {\LARGE\bfseries\textcolor{SZURed}{\coursename}}

    \vspace{0.4em}

    {\large \department}

    \vspace{1em}

    {\Large\bfseries\textcolor{SZUTextRed}{\homeworknumber}}

\end{center}

\vspace{1em}

% ============================================================
%                       学生信息
% ============================================================

\begin{center}

\begin{tabular}{rl@{\hspace{3cm}}rl}

\textcolor{SZUTextRed}{\bfseries 姓名：}
&
\studentname
&
\textcolor{SZUTextRed}{\bfseries 学号：}
&
\studentid

\end{tabular}

\end{center}

\vspace{0.7em}

{\color{SZUSoftBorder}\hrule}

\vspace{1.5em}



% ============================================================
% ============================================================
%
%                           题目 1
%
% ============================================================
% ============================================================

\begin{problem}[课本习题3]

设随机变量$(X,Y)$的概率密度为
\[
f(x,y)=
\begin{cases}
k(6-x-y), & 0<x<2,2<y<4,\\
0, & \text{其他}.
\end{cases}
\]
\begin{enumerate}[label=(\arabic*)]
    \item 确定常数$k$。
    \item 求$P\{X<1,Y<3\}$。
    \item 求$P\{X<1.5\}$。
    \item 求$P\{X+Y\le 4\}$。
\end{enumerate}
\end{problem}


% ============================================================
%
%              ↓↓↓↓↓↓ 学生答案区域 ↓↓↓↓↓↓
%
%          请只在 answer 环境中填写你的答案
%
% ============================================================

\begin{answer}

% ==========================================================
%                 请从下一行开始填写答案
% ==========================================================






% ==========================================================
%                 请在上一行结束填写答案
% ==========================================================

\end{answer}

% ============================================================
%
%              ↑↑↑↑↑↑ 学生答案区域 ↑↑↑↑↑↑
%
% ============================================================


\vspace{1em}



% ============================================================
% ============================================================
%
%                           题目 2
%
% ============================================================
% ============================================================

\begin{problem}[课本习题5]

设随机变量$(X,Y)$具有分布函数
\[
F(x,y)=
\begin{cases}
1-\mathrm{e}^{-x}-\mathrm{e}^{-y}+\mathrm{e}^{-x-y}, & x>0,y>0,\\
0, & \text{其他}.
\end{cases}
\]
求边缘分布函数。
\end{problem}


% ============================================================
%
%              ↓↓↓↓↓↓ 学生答案区域 ↓↓↓↓↓↓
%
% ============================================================

\begin{answer}

% ==========================================================
%                 请从下一行开始填写答案
% ==========================================================





% ==========================================================
%                 请在上一行结束填写答案
% ==========================================================

\end{answer}

% ============================================================
%
%              ↑↑↑↑↑↑ 学生答案区域 ↑↑↑↑↑↑
%
% ============================================================


\vspace{1em}



% ============================================================
% ============================================================
%
%                           题目 3
%
% ============================================================
% ============================================================

\begin{problem}[课本习题9]

设二维随机变量$(X,Y)$的概率密度为
\[
f(x,y)=
\begin{cases}
cx^2y, & x^2\le y\le 1,\\
0, & \text{其他}.
\end{cases}
\]
\begin{enumerate}[label=(\arabic*)]
    \item 确定常数$c$。
    \item 求边缘概率密度。
\end{enumerate}
\end{problem}


% ============================================================
%
%              ↓↓↓↓↓↓ 学生答案区域 ↓↓↓↓↓↓
%
% ============================================================

\begin{answer}

% ==========================================================
%                 请从下一行开始填写答案
% ==========================================================





% ==========================================================
%                 请在上一行结束填写答案
% ==========================================================

\end{answer}

% ============================================================
%
%              ↑↑↑↑↑↑ 学生答案区域 ↑↑↑↑↑↑
%
% ============================================================


\vspace{1em}



% ============================================================
% ============================================================
%
%                           题目 4
%
% ============================================================
% ============================================================

\begin{problem}[课本习题13]

在第9题中
\begin{enumerate}[label=(\arabic*)]
    \item 求条件概率密度$f_{X|Y}(x|y)$，特别，写出当$Y=\dfrac12$时$X$的条件概率密度。
    \item 求条件概率密度$f_{Y|X}(y|x)$，特别，分别写出当$X=\dfrac13,X=\dfrac12$时$Y$的条件概率密度。
    \item 求条件概率
    \[
    P\left\{Y\ge \frac14 \,\bigg|\, X=\frac12\right\},\quad P\left\{Y\ge \frac34 \,\bigg|\, X=\frac12\right\}.
    \]
\end{enumerate}
\end{problem}


% ============================================================
%
%              ↓↓↓↓↓↓ 学生答案区域 ↓↓↓↓↓↓
%
% ============================================================

\begin{answer}

% ==========================================================
%                 请从下一行开始填写答案
% ==========================================================





% ==========================================================
%                 请在上一行结束填写答案
% ==========================================================

\end{answer}

% ============================================================
%
%              ↑↑↑↑↑↑ 学生答案区域 ↑↑↑↑↑↑
%
% ============================================================


\vspace{1em}



% ============================================================
% ============================================================
%
%                           题目 5
%
% ============================================================
% ============================================================

\begin{problem}[课本习题15]

设随机变量$X\sim U(0,1)$，当给定$X=x$时，随机变量$Y$的条件概率密度为
\[
f_{Y|X}(y|x)=
\begin{cases}
x, & 0<y<\dfrac1x,\\
0, & \text{其他}.
\end{cases}
\]
\begin{enumerate}[label=(\arabic*)]
    \item 求$X$和$Y$的联合概率密度$f(x,y)$。
    \item 求边缘密度$f_Y(y)$，并画出它的图形。
    \item 求$P\{X>Y\}$。
\end{enumerate}
\end{problem}


% ============================================================
%
%              ↓↓↓↓↓↓ 学生答案区域 ↓↓↓↓↓↓
%
% ============================================================

\begin{answer}

% ==========================================================
%                 请从下一行开始填写答案
% ==========================================================





% ==========================================================
%                 请在上一行结束填写答案
% ==========================================================

\end{answer}

% ============================================================
%
%              ↑↑↑↑↑↑ 学生答案区域 ↑↑↑↑↑↑
%
% ============================================================


\vspace{1em}



% ============================================================
% ============================================================
%
%                           题目 6
%
% ============================================================
% ============================================================

\begin{problem}[课本习题17]

\begin{enumerate}[label=(\arabic*)]
    \item 设随机变量$(X,Y)$具有分布函数
    \[
    F(x,y)=
    \begin{cases}
    (1-\mathrm{e}^{-\alpha x})y, & x\ge0,0\le y\le1,\quad \alpha>0,\\
    1-\mathrm{e}^{-\alpha x}, & x\ge0,y>1,\\
    0, & \text{其他}.
    \end{cases}
    \]
    证明$X,Y$相互独立。
    \item 设随机变量$(X,Y)$具有分布律
    \[
    P\{X=x,Y=y\}=p^2(1-p)^{x+y-2},\quad 0<p<1,x,y\text{ 均为正整数},
    \]
    问$X,Y$是否相互独立。
\end{enumerate}
\end{problem}


% ============================================================
%
%              ↓↓↓↓↓↓ 学生答案区域 ↓↓↓↓↓↓
%
% ============================================================

\begin{answer}

% ==========================================================
%                 请从下一行开始填写答案
% ==========================================================





% ==========================================================
%                 请在上一行结束填写答案
% ==========================================================

\end{answer}

% ============================================================
%
%              ↑↑↑↑↑↑ 学生答案区域 ↑↑↑↑↑↑
%
% ============================================================





\end{document}
```
